% GENERATED FILE. Edit canonical lessons, then run npm run generate:books.
% canonical-source-hash: fnv1a64:dafb6da340bea71e
% canonical-lessons: IT-C305-alluvione, IT-C305-allagare, IT-C305-frana, IT-C305-valanga, IT-C305-annegare

\chapter{Acqua e fango}
\label{ch:it-acqua-e-fango}
\hlchaptermodality{305}

\begin{chapteropening}
\textbf{By the end of this chapter:} \emph{I can talk about floods, landslides and avalanches, and say that someone drowned.}

Complete the last lesson of chapter 305: I can talk about floods, landslides and avalanches, and say that someone drowned.
\end{chapteropening}

\section[l'alluvione]{l'alluvione — the flood}

\label{lesson:IT-C305-alluvione}



\begin{quote}
Before the new one: say the Italian for the evacuee, then the Italian for to evacuate.
\end{quote}

\subsection*{You'll want to know: l'alluvione}
\textbf{l'alluvione} — \textquotedblleft{}the flood\textquotedblright{}.

\begin{quote}
Dopo l'alluvione la piazza era piena di fango.
\end{quote}

\begin{quote}
\emph{After the flood the square was full of mud.}
\end{quote}

Feminine: \emph{un'alluvione}.

\begin{cousinweb}
\emph{alluvione} is a learned borrowing from Latin \emph{alluviō} \textquotedblleft{}an overflowing\textquotedblright{} (Lewis \& Short), from \emph{alluere} \textquotedblleft{}to wash against\textquotedblright{}.
\end{cousinweb}

\begin{grammarlens}[title={What you've built}]
One more word for telling, explaining and describing at B1.
\end{grammarlens}

\subsection*{Your turn}
\begin{itemize}
\raggedright
  \item \emph{Say it:} \emph{l'alluvione}
  \item \emph{Say it:} \emph{l'alluvione}, once more
  \item \emph{Say it:} a sentence of your own with \emph{l'alluvione}
  \item \emph{Recall:} say the Italian for the evacuee, then the Italian for to evacuate, then say \emph{l'alluvione} again
\end{itemize}

\begin{tcolorbox}[breakable,colback=teal!4,colframe=teal!35!black,title={Before you move on}]
What does l'alluvione mean? (\textquotedblleft{}The flood\textquotedblright{}.) Say it once more.
\end{tcolorbox}
\section[allagare]{allagare — to flood}

\label{lesson:IT-C305-allagare}



\begin{quote}
Before the new one: say the Italian for to evacuate, then the Italian for the flood.
\end{quote}

\subsection*{You'll want to know: allagare}
\textbf{allagare} — \textquotedblleft{}to flood\textquotedblright{}.

\begin{quote}
La pioggia ha allagato la cantina.
\end{quote}

\begin{quote}
\emph{The rain flooded the cellar.}
\end{quote}

\emph{Allagato} — flooded.

\begin{cousinweb}
\emph{allagare} is built on \emph{lago} \textquotedblleft{}lake\textquotedblright{}, from Latin \emph{lacus} (Lewis \& Short): to turn into a lake.
\end{cousinweb}

\begin{grammarlens}[title={What you've built}]
One more word for telling, explaining and describing at B1.
\end{grammarlens}

\subsection*{Your turn}
\begin{itemize}
\raggedright
  \item \emph{Say it:} \emph{allagare}
  \item \emph{Say it:} \emph{allagare}, once more
  \item \emph{Say it:} a sentence of your own with \emph{allagare}
  \item \emph{Recall:} say the Italian for to evacuate, then the Italian for the flood, then say \emph{allagare} again
\end{itemize}

\begin{tcolorbox}[breakable,colback=teal!4,colframe=teal!35!black,title={Before you move on}]
What does allagare mean? (\textquotedblleft{}To flood\textquotedblright{}.) Say it once more.
\end{tcolorbox}
\section[la frana]{la frana — the landslide}

\label{lesson:IT-C305-frana}



\begin{quote}
Before the new one: say the Italian for the flood, then the Italian for to flood.
\end{quote}

\subsection*{You'll want to know: la frana}
\textbf{la frana} — \textquotedblleft{}the landslide\textquotedblright{}.

\begin{quote}
Una frana ha chiuso la strada.
\end{quote}

\begin{quote}
\emph{A landslide has closed the road.}
\end{quote}

\emph{Essere una frana} — to be hopeless at something: \emph{sono una frana in matematica}.

\begin{cousinweb}
The origin of \emph{frana} is not settled; a link with Latin \emph{vorāgō} \textquotedblleft{}an abyss\textquotedblright{} has been proposed (Lewis \& Short gives \emph{vorāgō}).
\end{cousinweb}

\begin{grammarlens}[title={What you've built}]
One more word for telling, explaining and describing at B1.
\end{grammarlens}

\subsection*{Your turn}
\begin{itemize}
\raggedright
  \item \emph{Say it:} \emph{la frana}
  \item \emph{Say it:} \emph{la frana}, once more
  \item \emph{Say it:} a sentence of your own with \emph{la frana}
  \item \emph{Recall:} say the Italian for the flood, then the Italian for to flood, then say \emph{la frana} again
\end{itemize}

\begin{tcolorbox}[breakable,colback=teal!4,colframe=teal!35!black,title={Before you move on}]
What does la frana mean? (\textquotedblleft{}The landslide\textquotedblright{}.) Say it once more.
\end{tcolorbox}
\section[la valanga]{la valanga — the avalanche}

\label{lesson:IT-C305-valanga}



\begin{quote}
Before the new one: say the Italian for to flood, then the Italian for the landslide.
\end{quote}

\subsection*{You'll want to know: la valanga}
\textbf{la valanga} — \textquotedblleft{}the avalanche\textquotedblright{}.

\begin{quote}
Una valanga è caduta sulla strada.
\end{quote}

\begin{quote}
\emph{An avalanche came down on the road.}
\end{quote}

Also a flood of things: \emph{una valanga di lettere}.

\begin{cousinweb}
\emph{valanga} is usually traced to northern Italian dialects; the deeper origin is not settled (EtymDB has no entry).
\end{cousinweb}

\begin{grammarlens}[title={What you've built}]
One more word for telling, explaining and describing at B1.
\end{grammarlens}

\subsection*{Your turn}
\begin{itemize}
\raggedright
  \item \emph{Say it:} \emph{la valanga}
  \item \emph{Say it:} \emph{la valanga}, once more
  \item \emph{Say it:} a sentence of your own with \emph{la valanga}
  \item \emph{Recall:} say the Italian for to flood, then the Italian for the landslide, then say \emph{la valanga} again
\end{itemize}

\begin{tcolorbox}[breakable,colback=teal!4,colframe=teal!35!black,title={Before you move on}]
What does la valanga mean? (\textquotedblleft{}The avalanche\textquotedblright{}.) Say it once more.
\end{tcolorbox}
\section[annegare]{annegare — to drown}

\label{lesson:IT-C305-annegare}



\begin{quote}
Before the new one: say the Italian for the landslide, then the Italian for the avalanche.
\end{quote}

\subsection*{You'll want to know: annegare}
\textbf{annegare} — \textquotedblleft{}to drown\textquotedblright{}.

\begin{quote}
Un uomo è annegato nel fiume.
\end{quote}

\begin{quote}
\emph{A man drowned in the river.}
\end{quote}

With \emph{essere}: \emph{è annegato}.

\begin{cousinweb}
\emph{annegare} is inherited from Vulgar Latin \emph{adnecāre} (EtymDB), from Latin \emph{necāre} \textquotedblleft{}to kill\textquotedblright{}: first, to kill by water.
\end{cousinweb}

\begin{grammarlens}[title={What you've built}]
One more word for telling, explaining and describing at B1.
\end{grammarlens}

\subsection*{Your turn}
\begin{itemize}
\raggedright
  \item \emph{Say it:} \emph{annegare}
  \item \emph{Say it:} \emph{annegare}, once more
  \item \emph{Say it:} a sentence of your own with \emph{annegare}
  \item \emph{Recall:} say the Italian for the landslide, then the Italian for the avalanche, then say \emph{annegare} again
\end{itemize}

\begin{tcolorbox}[breakable,colback=teal!4,colframe=teal!35!black,title={Before you move on}]
What does annegare mean? (\textquotedblleft{}To drown\textquotedblright{}.) Say it once more.
\end{tcolorbox}
